\documentclass[../DoAn.tex]{subfiles}
\begin{document}

\lstdefinestyle{cmd}{
    basicstyle=\ttfamily\small,
    backgroundcolor=\color[gray]{0.95},
    frame=single,
    framerule=0.4pt,
    breaklines=true,
    breakatwhitespace=false,
    columns=flexible,
}

% ---------------------------------------------------------------
\section{Yêu cầu phần cứng và phần mềm}
\label{sec:pla-yeucau}

Hệ thống được phát triển và kiểm thử trên Windows 11 (64-bit) với hai cấu hình
máy tính khác nhau, liệt kê trong Bảng~\ref{tab:pla-hardware}.

\begin{table}[h!]
\centering
\caption{Cấu hình hai máy tính dùng để phát triển và kiểm thử}
\label{tab:pla-hardware}
\begin{tabular}{|l|l|l|}
\hline
\textbf{Thành phần} & \textbf{Máy bàn} & \textbf{Máy xách tay (Dell Inspiron 5620)} \\
\hline
CPU  & Intel Core i5-14400F (10 nhân) & Intel Core i5-1240P (12 nhân) \\
RAM  & 32 GB DDR4                     & 16 GB DDR4 \\
GPU  & NVIDIA RTX 4060                & NVIDIA MX570 \\
\hline
\end{tabular}
\end{table}

Phần mềm không đặt ra yêu cầu phần cứng tối thiểu cụ thể; hiệu năng
phụ thuộc chủ yếu vào số lượng phương tiện trong mô phỏng.
Cả hai cấu hình trên đều vận hành hệ thống ở chế độ thời gian thực với
kịch bản OSM thông thường mà không gặp vấn đề về hiệu năng.

\textbf{Phần mềm bắt buộc:}
\begin{itemize}
    \item \textbf{Python 3.10 trở lên} — ngôn ngữ triển khai server mô phỏng.
        Tải tại \url{https://www.python.org/downloads/}.
        Trong quá trình cài đặt, cần tích chọn \textit{``Add Python to PATH''}.
    \item \textbf{Eclipse SUMO} và thư viện \textbf{TraCI} — được cài tự động qua
        script \texttt{install.bat} (xem Mục~\ref{sec:pla-caidat}).
    \item \textbf{Microsoft Visual C++ Redistributable 2022 x64} — thư viện
        runtime cần thiết để chạy \texttt{sumo-gui.exe}.
        Cũng được cài tự động qua \texttt{install.bat}.
    \item \textbf{DirectX 11} trở lên — yêu cầu của bản build Unity.
        Đã có sẵn trên Windows 10/11; không cần cài thêm trên phần lớn máy tính hiện đại.
\end{itemize}

% ---------------------------------------------------------------
\section{Cài đặt môi trường}
\label{sec:pla-caidat}

Toàn bộ quá trình cài đặt được tự động hóa qua file \texttt{install.bat}
nằm trong thư mục gốc của gói phần mềm.
Script này thực hiện tuần tự các bước: kiểm tra Python, nâng cấp \texttt{pip},
cài gói \texttt{eclipse-sumo} (bao gồm SUMO, TraCI và sumolib), đồng thời
kiểm tra và cài Microsoft Visual C++ Redistributable nếu chưa có.

Để chạy script, nhấp đúp vào \texttt{install.bat} hoặc thực thi lệnh sau
trong Command Prompt tại thư mục gốc của gói phần mềm:

\begin{lstlisting}[style=cmd]
install.bat
\end{lstlisting}

Khi cài đặt thành công, script in ra phiên bản TraCI và đường dẫn
\texttt{SUMO\_HOME} như ví dụ dưới đây:

\begin{lstlisting}[style=cmd]
[OK] Python 3.12.0
[OK] Visual C++ Redistributable da co san.
[OK] traci 1.21.0
[OK] SUMO_HOME: C:\Users\...\site-packages\sumo
\end{lstlisting}

Để xác nhận thủ công, chạy lệnh sau trong terminal:

\begin{lstlisting}[style=cmd]
python -c "import traci; print(traci.VERSION)"
\end{lstlisting}

Nếu lệnh in ra số phiên bản mà không báo lỗi, môi trường đã sẵn sàng.

% ---------------------------------------------------------------
\section{Khởi chạy hệ thống}
\label{sec:pla-khoidong}

Hệ thống hỗ trợ hai chế độ vận hành: mô phỏng thời gian thực và phát lại
kịch bản đã ghi sẵn.

\subsection{Mô phỏng thời gian thực}
\label{subsec:pla-realtime}

Chế độ này khởi động SUMO song song với Unity, truyền dữ liệu vị trí phương
tiện theo từng bước mô phỏng. Thứ tự khởi chạy như sau.

\textbf{Bước 1 — Khởi động server mô phỏng.}
Nhấp đúp vào file \texttt{Simulation.bat}
(hoặc chạy trực tiếp trong terminal):

\begin{lstlisting}[style=cmd]
Simulation.bat
\end{lstlisting}

Giao diện launcher hiện ra với ba tab kịch bản:
\begin{itemize}
    \item \textbf{Mô phỏng Maze (.map)} — mô phỏng trên bản đồ mê cung
        được tạo sẵn từ file \texttt{.map}.
    \item \textbf{Mô phỏng OSM (.osm)} — mô phỏng trên bản đồ thực tế
        lấy từ OpenStreetMap; hệ thống tự động chuyển đổi file
        \texttt{.osm} sang định dạng SUMO.
    \item \textbf{Kịch bản netedit} — chạy kịch bản do người dùng
        tự xây dựng bằng netedit, chỉ cần trỏ đến thư mục chứa
        \texttt{.net.xml} và \texttt{.rou.xml}.
\end{itemize}
Sau khi chọn tab và cấu hình các tham số, nhấn nút \textbf{Khởi động Server}.
Launcher sẽ tự mở SUMO và bản build Unity.

\textbf{Bước 2 — Kết nối Unity và bắt đầu mô phỏng.}
Trong cửa sổ Unity, nhấn \textbf{Mô phỏng thời gian thực},
sau đó nhấn \textbf{Lấy dữ liệu đường} và chờ thông báo xác nhận
lấy dữ liệu thành công. Cuối cùng nhấn \textbf{Bắt đầu mô phỏng}
để bắt đầu quan sát.

\subsection{Phát lại kịch bản có sẵn (tiền kết xuất)}
\label{subsec:pla-prerender}

Chế độ tiền kết xuất không yêu cầu SUMO đang chạy — Unity đọc trực tiếp
từ hai file JSON đã ghi sẵn trong quá trình mô phỏng.

\textbf{Bước 1 — Chuẩn bị kịch bản.}
Một kịch bản hợp lệ gồm hai file:
\begin{itemize}
    \item \texttt{road\_data.json} — dữ liệu hình học mạng lưới đường.
    \item \texttt{scenario.json} — chuỗi trạng thái phương tiện theo từng bước.
\end{itemize}
Các file này được tạo tự động khi chạy mô phỏng thời gian thực và lưu vào
thư mục \texttt{Server/result/}.
Có thể dùng lại các kịch bản đã ghi trước đó.

\textbf{Bước 2 — Mở bản build Unity và cấu hình đường dẫn.}
Mở file \texttt{TestGR1.1.exe} trong thư mục \texttt{UnityBuild/},
chọn \textbf{Tiền kết xuất}, nhập đường dẫn tuyệt đối của
\texttt{road\_data.json} vào ô \textit{``Đường dẫn bản đồ''} và
đường dẫn tuyệt đối của \texttt{scenario.json} vào ô
\textit{``Đường dẫn kịch bản''}.

\textbf{Bước 3 — Phát lại.}
Nhấn nút \textbf{Phát lại}. Unity sẽ tải dữ liệu và bắt đầu tái hiện
kịch bản mà không cần kết nối mạng hay SUMO.

% ---------------------------------------------------------------
\section{Xử lý sự cố thường gặp}
\label{sec:pla-sucố}

\textbf{Lỗi kết nối TCP (``Connection refused'' hoặc ``Connection timed out'').}
Server Python và Unity giao tiếp qua các cổng TCP 5050, 5053 và 5054.
Nếu gặp lỗi kết nối, cần kiểm tra: (i) launcher đã được khởi động trước
Unity; (ii) tường lửa Windows không chặn các cổng trên;
(iii) không có tiến trình SUMO hoặc Python nào từ phiên trước còn đang giữ cổng
(có thể kết thúc qua Task Manager).

\textbf{\texttt{sumo-gui} báo lỗi thiếu DLL.}
Nguyên nhân thường là thiếu Microsoft Visual C++ Redistributable.
Chạy lại \texttt{install.bat} để cài đặt tự động, hoặc tải thủ công tại
\url{https://aka.ms/vs/17/release/vc_redist.x64.exe}.

\textbf{\texttt{import traci} báo \texttt{ModuleNotFoundError}.}
Chạy lại \texttt{install.bat} để cài lại gói \texttt{eclipse-sumo}.
Đảm bảo lệnh \texttt{python} trong terminal trỏ đúng vào bản Python đã cài
gói (kiểm tra bằng \texttt{where python}).

\end{document}
